\section{Bounds on the smallest eigenvalue of regular graphs} 

Let $-1+\lambda$ denote the smallest eigenvalue of a $d$-regular, non-bipartite graph $\Gamma$. Since~$\Gamma$ is non-bipartite, it follows that $\lambda \neq 0$. So, the second largest eigenvalue of the normalised adjacency operator of $\Gamma^2$ is $\geq (1- \lambda)^2$. This implies that the smallest nontrivial eigenvalue of the normalised Laplacian operator of $\Gamma^2$ is $\leq 1 - (1-\lambda)^2$. By the discrete Cheeger--Buser inequality~\cite[Theorem 2.2]{ChungSpectralGraphTheory}, it follows that the edge-Cheeger constant of $\Gamma^2$ is $\leq \sqrt {2 (1 - (1-\lambda)^2)} < 2\sqrt \lambda$. Hence, if the vertex-Cheeger constant of $\Gamma^2$ is $\geq \varpi$, then the inequality $2 \sqrt \lambda > \frac{\varpi}{d^2}$ holds, and hence the smallest eigenvalue of $\Gamma$ is $> -1 + \frac {\varpi^2}{4d^4}$. 
The above discussion, combined with Theorem~\ref{Thm:SpannSubgraVertexTra}, yields the following. 

\begin{proposition}
\label{Cor:AlmostGraphAut}
Suppose $(V, E)$ is regular and it is nearly vertex-transitive under~$\calG$. Then the smallest eigenvalue of the normalised adjacency operator of $(V, E)$ is greater than $-1+ \frac {\varepsilon^4}{2^{8}3^{4}d^6}$. 
Assume in addition that $(V, E)$ is nearly vertex-transitive under $\calG$ with respect to $\{\psi_{i,v}\}$ and $\psi_{i, v}$ sends any index two subgroup of $\calG$ to itself. Then the smallest eigenvalue of the normalised adjacency operator of $(V, E)$ is greater than $ -1+\frac {\varepsilon^4}{60^2 d^6}$. 
\end{proposition}

\begin{proposition}
\label{Cor:SpannSubgraVertexTra}
If $(V, E)$ is vertex-transitive, then the smallest eigenvalue of its normalised adjacency operator is greater than $ -1+\frac {\varepsilon^4}{60^2 d^{10}}$. 
\end{proposition}
